% Generated draft: requires mathematical and source review.
\documentclass[11pt]{article}
\usepackage[margin=25mm]{geometry}
\usepackage{fontspec}
\setmainfont{Latin Modern Roman}
\usepackage{amsmath,amssymb,mathtools,unicode-math}
\setmathfont{Latin Modern Math}
\usepackage{enumitem,needspace,xurl,hyperref}
\setlength{\parindent}{0pt}
\setlength{\parskip}{5pt}
\setlength{\emergencystretch}{3em}
\newcommand{\N}{\mathbb{N}}
\newcommand{\Z}{\mathbb{Z}}
\newcommand{\Q}{\mathbb{Q}}
\newcommand{\R}{\mathbb{R}}
\newcommand{\C}{\mathbb{C}}
\newcommand{\F}{\mathbb{F}}
\newcommand{\A}{\mathbb{A}}
\newcommand{\PP}{\mathbb{P}}
\newcommand{\E}{\mathbb{E}}
\newcommand{\eps}{\varepsilon}
\newcommand{\dd}{\,\mathrm{d}}
\begin{document}
\section*{UnsolvedMath ID: 30004310}
\textbf{Source problem code:} catalogue entry
\subsection{Definitions and mathematical statement}
For a finite nonempty set \(X\), write \(X^2=X\times X\). A map \(r:X^2\to X^2\) satisfies the set-theoretic Yang--Baxter equation if
\[
(r\times\mathrm{id}_X)(\mathrm{id}_X\times r)(r\times\mathrm{id}_X)
=
(\mathrm{id}_X\times r)(r\times\mathrm{id}_X)(\mathrm{id}_X\times r)
\quad\text{on }X^3.
\]
Its size is \(|X|\). Write \(r(x,y)=(\sigma_x(y),\tau_y(x))\). Nondegeneracy means that every \(\sigma_x:X\to X\) and every \(\tau_y:X\to X\) is bijective; involutivity means \(r^2=\mathrm{id}_{X^2}\).

For nondegenerate involutive set-theoretic solutions, the usual retraction identifies \(x,y\in X\) when \(\sigma_x=\sigma_y\); the quotient carries the map induced by \(r\). With \(\operatorname{Ret}^0(X,r)=(X,r)\) and subsequent retractions defined iteratively, such a solution is called a multipermutation solution if \(\operatorname{Ret}^m(X,r)\) has one element for some integer \(m\geq 0\). This definition specifies the usual nondegenerate involutive convention only; the supplied record does not establish that this is the intended class.

To express the counting target, fix a class \(\mathcal C\) of finite solutions, a definition of multipermutation within that class, and a counting convention. For labelled counting, let \(\mathcal C_n\) consist of the eligible maps on \(\{1,\ldots,n\}\). For counting up to isomorphism, let \(\mathcal C_n\) instead consist of their isomorphism classes, where \((X,r)\) and \((Y,s)\) are isomorphic if a bijection \(f:X\to Y\) satisfies
\[
(f\times f)\circ r=s\circ(f\times f).
\]
Let \(\mathcal M_n\subseteq\mathcal C_n\) be the multipermutation solutions under the same convention, and put \(a_n=|\mathcal C_n|\), \(b_n=|\mathcal M_n|\).

\textbf{Question, conditional on these conventions.} Provided \(a_n>0\) for all sufficiently large positive integers \(n\), is
\[
\lim_{n\to\infty}\frac{b_n}{a_n}=1?
\]
Equivalently, does every \(\varepsilon>0\) admit an integer \(N\geq1\) such that, for every integer \(n\geq N\), \(a_n>0\) and \(b_n\geq(1-\varepsilon)a_n\)?

The source does not specify the eligible class, including whether bijectivity, nondegeneracy, involutivity, square-freeness, or other restrictions are required, nor whether counting is labelled or up to isomorphism. Consequently this is a conditional formulation of the supplied asymptotic question, pending clarification of those choices.
\subsection{Short English statement}
The question asks whether multipermutation solutions constitute an asymptotic proportion of \(1\) among Yang--Baxter solutions of size \(n\). The eligible solutions and counting convention are unspecified, so the precise counting problem requires clarification. The supplied record labels it open; that status is not independently verified here.
\subsection{Sources}
\begin{itemize}
\item Supplied local UnsolvedMath registry record, permanent ID \(30004310\), title \emph{Prevalence of Multipermutation Yang--Baxter Solutions}, source problem code \texttt{catalogue entry}, in the Oberwolfach Reports Open Problems collection. The proposed year is recorded as \(2019\); no individual proposer is supplied.
\item Supplied report citation: \emph{Mini-Workshop: Algebraic Tools for Solving the Yang--Baxter Equation (2020)}, Oberwolfach Reports, pp.~3207--3242, DOI \texttt{10.4171/owr/2019/51}. \url{https://doi.org/10.4171/owr/2019/51}. The background describes Oberwolfach Reports as documenting mathematical workshops and their open-problem discussions.
\item Supplied literature source: \emph{Incremental SAT-based enumeration of Yang--Baxter solutions}. \url{https://link.springer.com/chapter/10.1007/978-3-031-90653-4_1}.
\item Supplied literature source: Ferrara, Trombetti, and Tsang, \emph{Cardinalities and multipermutation level}. \url{https://arxiv.org/abs/2509.11001}.
\item Supplied literature source: Castelli and Trappeniers, \emph{Finite indecomposable involutive solutions}. \url{https://arxiv.org/abs/2303.00581}.
\item Literature assessment date, copied exactly: \texttt{2026-08-22}. The supplied dated curation labels the problem open and states that finite enumeration and structural results for multipermutation solutions do not establish the proposed asymptotic ratio. The background characterizes this as dated literature triage derived from the supplied OWR curation and requests independent verification of time-sensitive status and cited claims. This assessment is supplied curation, not independently verified here.
\item The registry gives the statement status as \texttt{corrected\_verified} and reports a manual semantic, typographic, and LaTeX review asserting that the clean statement is self-contained and preserves the target without unsupported additions. These are supplied registry assertions; they do not resolve the formulation ambiguities identified in this draft.
\end{itemize}
\end{document}
